%% ============================================================
%% cap05.tex — Capítulo 5: Implementación en Rust
%% Programación Orientada a Objetos: Diseño e Implementación
%% de un Sistema de Información Académico-Administrativo
%% para la Universidad Distrital Francisco José de Caldas
%%
%% Autor: Roberto Albeiro Pava-Díaz, Ph.D.
%% ============================================================

\chapter{Implementación en \rust}
\label{cap:rust}


\noindent
El estudio de \rust\ en un libro de \poo\ requiere una aclaración inicial, toda
vez que \rust\ no es un lenguaje orientado a objetos en el sentido clásico. No
posee clases, no admite herencia de structs y no tiene un modelo de objetos con
identidad mutable implícita. Sin embargo, los principios fundamentales de la
orientación a objetos ---encapsulamiento, abstracción, polimorfismo--- se logran
en \rust\ mediante mecanismos disímiles: los \texttt{struct} encapsulan estado,
los \texttt{trait} definen comportamiento abstracto, y el \textit{despacho dinámico}
(\texttt{dyn Trait}) habilita el polimorfismo en tiempo de ejecución.\footnote{Esta
aproximación pragmática a la POO en Rust ha sido caracterizada por \textcite{Klabnik2023}
como ``composición sobre herencia'', principio que coincide con las recomendaciones
de \textcite{Gamma1994} para el diseño orientado a objetos flexible.}

El propósito de implementar el mismo dominio estatutario en \rust\ no es demostrar
que un lenguaje es superior al otro. Es suscitar reflexiones sobre por qué ciertos
conceptos de la POO ---herencia, excepciones, referencias nulas--- se diseñaron de
determinada manera y qué costos y beneficios conllevan. Queda la duda ---y
deliberadamente se deja abierta--- de si el modelo de tipos de \rust\ habría evitado
errores históricos en sistemas de información universitaria cuya falla se debió,
precisamente, a referencias nulas no controladas.

% ─────────────────────────────────────────────────────────────────────────────
\section{Fundamentos de \rust}
\label{sec:rust-fundamentos}
% ─────────────────────────────────────────────────────────────────────────────

\noindent
Antes de examinar la implementación del dominio, conviene establecer los tres pilares
conceptuales de \rust\ que mayor impacto tienen en el diseño del sistema: el modelo
de propiedad (\textit{ownership}), el préstamo de referencias (\textit{borrowing})
con sus tiempos de vida (\textit{lifetimes}), y el sistema de tipos que reemplaza
los valores nulos y las excepciones.

\subsection{Ownership, borrowing y lifetimes}
\label{subsec:ownership}

El sistema de propiedad de \rust\ establece que cada valor tiene un único propietario
(\textit{owner}) en un instante dado. Cuando el propietario sale del alcance
(\textit{scope}), el valor es liberado automáticamente, sin recolector de basura.
Este mecanismo elimina las fugas de memoria por diseño, no por convención.\footnote{El
modelo de ownership fue formalizado por Matsakis y Klock (\cite{matsakis2014}) y
constituye la contribución más original de Rust al diseño de lenguajes de sistemas.}

\begin{lstlisting}[style=rustStyle, caption={Ownership y borrowing aplicados al dominio}, label={lst:ownership}]
// Demostración de ownership y borrowing
// usando las entidades del dominio estatutario

/// Representa una Escuela con su lista de docentes.
/// La Escuela es propietaria (owner) de sus docentes.
struct Escuela {
    nombre:   String,
    codigo:   String,
    docentes: Vec<Docente>,
}

struct Docente {
    nombre:         String,
    identificacion: String,
    tipo:           TipoDocente,
}

// ── Ejemplo 1: Move semantics ──────────────────────────────────────────────

fn ejemplo_move() {
    let escuela1 = Escuela {
        nombre:   String::from("Ingeniería de Sistemas"),
        codigo:   String::from("EIS"),
        docentes: Vec::new(),
    };
    // escuela1 se mueve (move) a escuela2; ya no es accesible por escuela1
    let escuela2 = escuela1;
    println!("Escuela: {}", escuela2.nombre);
    // println!("{}", escuela1.nombre); // ERROR de compilación: moved value
}

// ── Ejemplo 2: Borrowing inmutable ─────────────────────────────────────────

/// Imprime el nombre de la escuela sin tomar propiedad.
/// El parámetro es una referencia inmutable (&Escuela).
fn imprimir_escuela(escuela: &Escuela) {
    println!("Escuela: {} ({})", escuela.nombre, escuela.codigo);
    // No se puede modificar escuela aquí
}

// ── Ejemplo 3: Borrowing mutable ───────────────────────────────────────────

/// Agrega un docente a la escuela.
/// El parámetro es una referencia mutable (&mut Escuela).
/// En un instante dado solo puede existir UNA referencia mutable.
fn agregar_docente(escuela: &mut Escuela, docente: Docente) {
    escuela.docentes.push(docente);
}

// ── Ejemplo 4: Lifetimes ────────────────────────────────────────────────────

/// Retorna una referencia al docente con mayor antigüedad.
/// El lifetime 'a garantiza que la referencia retornada
/// no vive más que la escuela de donde proviene.
fn docente_mas_antiguo<'a>(escuela: &'a Escuela) -> Option<&'a Docente> {
    escuela.docentes.first()
}

fn main() {
    let mut escuela = Escuela {
        nombre:   String::from("Ing. de Sistemas"),
        codigo:   String::from("EIS"),
        docentes: Vec::new(),
    };

    let docente = Docente {
        nombre:         String::from("Prof. Ejemplo"),
        identificacion: String::from("12345678"),
        tipo:           TipoDocente::Planta,
    };

    agregar_docente(&mut escuela, docente);
    imprimir_escuela(&escuela);

    if let Some(d) = docente_mas_antiguo(&escuela) {
        println!("Primer docente: {}", d.nombre);
    }
}
\end{lstlisting}

\noindent
El listado~\ref{lst:ownership} ilustra un aspecto que los estudiantes provenientes
de \java\ encuentran sorprendente: la operación \texttt{let escuela2 = escuela1}
no copia el objeto, sino que transfiere la propiedad (\textit{move}). La variable
\texttt{escuela1} queda inválida después de esa línea. Este comportamiento, que
inicialmente parece una restricción arbitraria, es la garantía de que nunca
existirán dos propietarios simultáneos del mismo dato: condición suficiente para
eliminar las condiciones de carrera (\textit{data races}) en programas concurrentes.

\subsection{Structs, enums y traits}
\label{subsec:structs-traits}

Los \texttt{struct} de \rust\ encapsulan datos; los \texttt{trait} definen el
comportamiento abstracto que los structs pueden implementar. Esta separación entre
datos y comportamiento difiere del modelo de clases de \java\ ---donde datos y métodos
coexisten en la misma unidad--- pero permite una composición más flexible, puesto que
un struct puede implementar múltiples traits sin las restricciones del grafo de herencia.

\begin{lstlisting}[style=rustStyle, caption={Structs, enums y traits del dominio}, label={lst:structs-traits}]
use std::fmt;

// ─── Enumeraciones ──────────────────────────────────────────────────────────

/// Tipos de vinculación docente (Acuerdo 004 de 2025, Art. 8).
#[derive(Debug, Clone, PartialEq)]
pub enum TipoDocente {
    Planta,
    Ocasional,
    HoraCatedra,
    Visitante,
    Experto,
}

impl fmt::Display for TipoDocente {
    fn fmt(&self, f: &mut fmt::Formatter<'_>) -> fmt::Result {
        match self {
            TipoDocente::Planta      => write!(f, "Planta"),
            TipoDocente::Ocasional   => write!(f, "Ocasional"),
            TipoDocente::HoraCatedra => write!(f, "Hora Cátedra"),
            TipoDocente::Visitante   => write!(f, "Visitante"),
            TipoDocente::Experto     => write!(f, "Experto"),
        }
    }
}

/// Estados del ciclo de vida estudiantil.
#[derive(Debug, Clone, PartialEq)]
pub enum EstadoEstudiante {
    Aspirante,
    Admitido,
    Activo,
    Suspendido,
    Graduado,
    Retirado,
}

// ─── Trait: comportamiento de entidad académica ─────────────────────────────

/// Trait que define el comportamiento común de toda entidad académica.
/// Equivale a la clase abstracta UnidadAcademica de Java, pero sin herencia.
pub trait EntidadAcademica {
    fn nombre(&self) -> &str;
    fn codigo(&self) -> &str;
    fn tipo(&self)   -> &str;

    /// Implementación por defecto: descripción resumida de la entidad.
    fn descripcion(&self) -> String {
        format!("{} [{}]: {}", self.tipo(), self.codigo(), self.nombre())
    }
}

// ─── Structs organizacionales ───────────────────────────────────────────────

/// Facultad universitaria de primer nivel.
#[derive(Debug, Clone)]
pub struct Facultad {
    pub nombre:   String,
    pub codigo:   String,
    pub decano:   String,
    pub escuelas: Vec<Escuela>,
}

impl EntidadAcademica for Facultad {
    fn nombre(&self) -> &str { &self.nombre }
    fn codigo(&self) -> &str { &self.codigo }
    fn tipo(&self)   -> &str { "Facultad"  }
}

impl Facultad {
    pub fn new(nombre: &str, codigo: &str, decano: &str) -> Self {
        Facultad {
            nombre:   nombre.to_string(),
            codigo:   codigo.to_string(),
            decano:   decano.to_string(),
            escuelas: Vec::new(),
        }
    }

    pub fn agregar_escuela(&mut self, escuela: Escuela) {
        self.escuelas.push(escuela);
    }

    pub fn escuelas(&self) -> &[Escuela] {
        &self.escuelas
    }
}

// ─── Escuela ────────────────────────────────────────────────────────────────

#[derive(Debug, Clone)]
pub struct Escuela {
    pub nombre:   String,
    pub codigo:   String,
    pub director: String,
    pub docentes: Vec<Docente>,
}

impl EntidadAcademica for Escuela {
    fn nombre(&self) -> &str { &self.nombre }
    fn codigo(&self) -> &str { &self.codigo }
    fn tipo(&self)   -> &str { "Escuela"   }
}

impl Escuela {
    pub fn new(nombre: &str, codigo: &str, director: &str) -> Self {
        Escuela {
            nombre:   nombre.to_string(),
            codigo:   codigo.to_string(),
            director: director.to_string(),
            docentes: Vec::new(),
        }
    }

    pub fn agregar_docente(&mut self, d: Docente) {
        self.docentes.push(d);
    }
}

// ─── Centro e Instituto ─────────────────────────────────────────────────────

#[derive(Debug, Clone)]
pub struct Centro {
    pub nombre:   String,
    pub codigo:   String,
    pub director: String,
}

impl EntidadAcademica for Centro {
    fn nombre(&self) -> &str { &self.nombre }
    fn codigo(&self) -> &str { &self.codigo }
    fn tipo(&self)   -> &str { "Centro"    }
}

#[derive(Debug, Clone)]
pub struct Instituto {
    pub nombre:   String,
    pub codigo:   String,
    pub director: String,
}

impl EntidadAcademica for Instituto {
    fn nombre(&self) -> &str { &self.nombre }
    fn codigo(&self) -> &str { &self.codigo }
    fn tipo(&self)   -> &str { "Instituto" }
}

/// Comité Académico Básico de Área.
#[derive(Debug, Clone)]
pub struct Caba {
    pub area_formacion: String,
    pub coordinador:    String,
    pub nombre:         String,
}
\end{lstlisting}

\subsection{Pattern matching con \texttt{Result} y \texttt{Option}}
\label{subsec:pattern-matching}

El \textit{pattern matching} de \rust\ constituye su mecanismo más expresivo para
el manejo de valores que pueden estar ausentes (\texttt{Option<T>}) o ser el resultado
de operaciones que pueden fallar (\texttt{Result<T, E>}). En lugar de las excepciones
y los valores nulos de \java, \rust\ hace explícita en el tipo de retorno la posibilidad
de ausencia o error, constriñendo al llamador a tratar ambos casos en tiempo de compilación.

\begin{lstlisting}[style=rustStyle, caption={Pattern matching con \texttt{Option} y \texttt{Result}}, label={lst:pattern-matching}]
// ─── Option: búsqueda de estudiante por código ──────────────────────────────

pub struct Estudiante {
    pub nombre:    String,
    pub codigo:    String,
    pub programa:  String,
    pub semestre:  u8,
    pub estado:    EstadoEstudiante,
}

/// Busca un estudiante por código en una lista.
/// Retorna Option<&Estudiante>: Some si existe, None si no.
pub fn buscar_estudiante<'a>(
    estudiantes: &'a [Estudiante],
    codigo: &str,
) -> Option<&'a Estudiante> {
    estudiantes.iter().find(|e| e.codigo == codigo)
}

fn ejemplo_option() {
    let estudiantes = vec![
        Estudiante {
            nombre:   "Ana Pérez".to_string(),
            codigo:   "20201015".to_string(),
            programa: "ISC".to_string(),
            semestre: 6,
            estado:   EstadoEstudiante::Activo,
        },
    ];

    // Caso 1: estudiante encontrado
    match buscar_estudiante(&estudiantes, "20201015") {
        Some(e) => println!("Encontrado: {} — semestre {}", e.nombre, e.semestre),
        None    => println!("Estudiante no encontrado."),
    }

    // Caso 2: estudiante no encontrado
    // El compilador obliga a tratar AMBOS casos.
    // No es posible ignorar None sin advertencia.
    let resultado = buscar_estudiante(&estudiantes, "99999999");
    let nombre = resultado
        .map(|e| e.nombre.as_str())
        .unwrap_or("Desconocido");
    println!("Nombre: {}", nombre);
}

// ─── Result: validación de matrícula ────────────────────────────────────────

#[derive(Debug)]
pub enum ErrorMatricula {
    ProgramaInexistente(String),
    EstudianteSuspendido,
    PeriodoFuera,
}

impl fmt::Display for ErrorMatricula {
    fn fmt(&self, f: &mut fmt::Formatter<'_>) -> fmt::Result {
        match self {
            ErrorMatricula::ProgramaInexistente(cod) =>
                write!(f, "Programa '{}' no existe en el sistema.", cod),
            ErrorMatricula::EstudianteSuspendido =>
                write!(f, "El estudiante se encuentra suspendido."),
            ErrorMatricula::PeriodoFuera =>
                write!(f, "El periodo de matrícula está cerrado."),
        }
    }
}

/// Valida y registra una matrícula. Retorna Result.
pub fn matricular(
    estudiante: &Estudiante,
    programa_codigo: &str,
    periodo_abierto: bool,
) -> Result<String, ErrorMatricula> {
    if !periodo_abierto {
        return Err(ErrorMatricula::PeriodoFuera);
    }
    if estudiante.estado == EstadoEstudiante::Suspendido {
        return Err(ErrorMatricula::EstudianteSuspendido);
    }
    // En un sistema real se consultaría el repositorio de programas
    if programa_codigo.is_empty() {
        return Err(ErrorMatricula::ProgramaInexistente(
            programa_codigo.to_string()));
    }
    Ok(format!("Matrícula registrada: {} en {} — periodo activo",
               estudiante.nombre, programa_codigo))
}

fn ejemplo_result() {
    let e = Estudiante {
        nombre:   "Carlos Gómez".to_string(),
        codigo:   "20191234".to_string(),
        programa: "ISC".to_string(),
        semestre: 4,
        estado:   EstadoEstudiante::Activo,
    };

    // El operador ? propaga el error automáticamente
    match matricular(&e, "ISC-2025", true) {
        Ok(msg) => println!("{}", msg),
        Err(err) => println!("Error de matrícula: {}", err),
    }
}
\end{lstlisting}

% ─────────────────────────────────────────────────────────────────────────────
\section{Modelado del dominio completo en \rust}
\label{sec:rust-dominio}
% ─────────────────────────────────────────────────────────────────────────────

\noindent
La implementación completa del dominio en \rust\ abarca los actores institucionales,
los programas académicos y el sistema de evaluación. En todos los casos se prioriza
la seguridad de tipos y la expresividad del sistema de pattern matching frente a la
herencia de clases que caracteriza la versión en \java.

\subsection{Structs y enums de actores}
\label{subsec:rust-actores}

\begin{lstlisting}[style=rustStyle, caption={Actores del dominio en \rust}, label={lst:rust-actores}]
use serde::{Deserialize, Serialize};

// ─── Docente ────────────────────────────────────────────────────────────────

#[derive(Debug, Clone, Serialize, Deserialize)]
pub struct Docente {
    pub nombre:         String,
    pub identificacion: String,
    pub correo:         Option<String>,  // puede no tener correo registrado
    pub tipo:           TipoDocente,
    pub escuela_codigo: String,
    pub categoria:      Option<String>,  // sólo para planta
}

impl Docente {
    pub fn new(
        nombre: &str,
        identificacion: &str,
        tipo: TipoDocente,
        escuela_codigo: &str,
    ) -> Self {
        Docente {
            nombre:         nombre.to_string(),
            identificacion: identificacion.to_string(),
            correo:         None,
            tipo,
            escuela_codigo: escuela_codigo.to_string(),
            categoria:      None,
        }
    }

    pub fn con_correo(mut self, correo: &str) -> Self {
        self.correo = Some(correo.to_string());
        self
    }

    pub fn con_categoria(mut self, categoria: &str) -> Self {
        self.categoria = Some(categoria.to_string());
        self
    }

    /// ¿Es docente de planta? Verificación por pattern matching.
    pub fn es_planta(&self) -> bool {
        matches!(self.tipo, TipoDocente::Planta)
    }
}

// ─── PersonalAdministrativo ─────────────────────────────────────────────────

#[derive(Debug, Clone, Serialize, Deserialize)]
pub struct PersonalAdministrativo {
    pub nombre:         String,
    pub identificacion: String,
    pub correo:         Option<String>,
    pub cargo:          String,
    pub dependencia:    String,
}

// ─── Egresado ───────────────────────────────────────────────────────────────

#[derive(Debug, Clone, Serialize, Deserialize)]
pub struct Egresado {
    pub nombre:               String,
    pub identificacion:       String,
    pub programa_graduacion:  String,
    pub anio_graduacion:      u16,
    pub titulo:               String,
}

// ─── Programa Académico ─────────────────────────────────────────────────────

#[derive(Debug, Clone, PartialEq, Serialize, Deserialize)]
pub enum NivelPrograma {
    Pregrado,
    Especializacion,
    Maestria,
    Doctorado,
}

#[derive(Debug, Clone, Serialize, Deserialize)]
pub struct ProgramaAcademico {
    pub nombre:          String,
    pub codigo:          String,
    pub nivel:           NivelPrograma,
    pub modalidad:       String,
    pub facultad_codigo: String,
    pub plan_estudios:   Option<PlanDeEstudios>,
}

#[derive(Debug, Clone, Serialize, Deserialize)]
pub struct PlanDeEstudios {
    pub version:          String,
    pub creditos_totales: u16,
    pub espacios:         Vec<EspacioAcademico>,
}

impl PlanDeEstudios {
    pub fn creditos_acumulados(&self) -> u16 {
        self.espacios.iter().map(|e| e.creditos).sum()
    }
}

#[derive(Debug, Clone, Serialize, Deserialize)]
pub struct EspacioAcademico {
    pub nombre:         String,
    pub codigo:         String,
    pub creditos:       u16,
    pub tipo:           String,
    pub horas_teoricas: u8,
    pub horas_practicas: u8,
}

// ─── Matrícula ──────────────────────────────────────────────────────────────

#[derive(Debug, Clone, Serialize, Deserialize)]
pub struct Matricula {
    pub estudiante_id:   String,
    pub programa_codigo: String,
    pub periodo:         String,
    pub estado:          EstadoEstudiante,
}
\end{lstlisting}

\subsection{Clases de evaluación en \rust}
\label{subsec:rust-evaluacion}

\begin{lstlisting}[style=rustStyle, caption={Sistema de evaluación en \rust}, label={lst:rust-evaluacion}]
use std::collections::HashMap;

// ─── Resultado de Aprendizaje ────────────────────────────────────────────────

#[derive(Debug, Clone, Serialize, Deserialize)]
pub struct ResultadoDeAprendizaje {
    pub descripcion:     String,
    pub nivel_dominio:   String,
    pub programa_codigo: String,
}

// ─── Rúbrica ────────────────────────────────────────────────────────────────

#[derive(Debug, Clone, Serialize, Deserialize)]
pub struct Rubrica {
    pub nombre:    String,
    pub criterios: Vec<CriterioRubrica>,
    pub niveles:   Vec<String>,
}

#[derive(Debug, Clone, Serialize, Deserialize)]
pub struct CriterioRubrica {
    pub nombre:       String,
    pub ponderacion:  f64,
    pub descriptores: HashMap<String, String>,
}

impl CriterioRubrica {
    pub fn nuevo(nombre: &str, ponderacion: f64) -> Result<Self, String> {
        if !(0.0..=1.0).contains(&ponderacion) {
            return Err(format!(
                "Ponderación {ponderacion} fuera de rango [0.0, 1.0]"));
        }
        Ok(CriterioRubrica {
            nombre:       nombre.to_string(),
            ponderacion,
            descriptores: HashMap::new(),
        })
    }
}

// ─── Evidencia ──────────────────────────────────────────────────────────────

#[derive(Debug, Clone, Serialize, Deserialize)]
pub struct Evidencia {
    pub tipo:          String,
    pub descripcion:   String,
    pub fecha:         String,  // ISO 8601: "2025-11-15"
    pub calificacion:  f64,     // 0.0 - 5.0
}

impl Evidencia {
    pub fn nueva(
        tipo: &str,
        descripcion: &str,
        fecha: &str,
        calificacion: f64,
    ) -> Result<Self, String> {
        if !(0.0..=5.0).contains(&calificacion) {
            return Err(format!(
                "Calificación {calificacion} fuera del rango 0.0–5.0."));
        }
        Ok(Evidencia {
            tipo:         tipo.to_string(),
            descripcion:  descripcion.to_string(),
            fecha:        fecha.to_string(),
            calificacion,
        })
    }

    pub fn aprueba(&self, umbral: f64) -> bool {
        self.calificacion >= umbral
    }
}

// ─── Trait Evaluable ────────────────────────────────────────────────────────

/// Trait para entidades sujetas a evaluación formativa.
pub trait Evaluable {
    fn evaluar(&self, rubrica: &Rubrica) -> Result<f64, String>;
    fn evidencias(&self) -> &[Evidencia];

    fn aprueba(&self, umbral: f64) -> bool {
        self.evaluar(&Rubrica {
            nombre: String::new(),
            criterios: Vec::new(),
            niveles: Vec::new(),
        }).map(|nota| nota >= umbral)
          .unwrap_or(false)
    }
}

// ─── Trait Gestionable ──────────────────────────────────────────────────────

/// Trait para operaciones CRUD sobre entidades del dominio.
pub trait Gestionable {
    fn crear(&mut self) -> Result<(), String>;
    fn actualizar(&mut self) -> Result<(), String>;
    fn eliminar(&mut self) -> Result<(), String>;
}
\end{lstlisting}

% ─────────────────────────────────────────────────────────────────────────────
\section{Concurrencia y seguridad de memoria}
\label{sec:rust-concurrencia}
% ─────────────────────────────────────────────────────────────────────────────

\noindent
Uno de los aportes más significativos de \rust\ es la garantía en tiempo de compilación
de la ausencia de condiciones de carrera (\textit{data races}). Este aspecto resulta
relevante para el sistema estatutario en escenarios como el procesamiento concurrente
de matrículas durante los periodos de inscripción, cuando múltiples usuarios intentan
registrar matrícula en el mismo programa simultáneamente.

\begin{lstlisting}[style=rustStyle, caption={caption[Acceso concurrente con Arc y Mutex]{Acceso concurrente con \texttt{Arc$<$Mutex$>$}}}, label={lst:concurrencia}]
use std::sync::{Arc, Mutex};
use std::thread;

/// Procesamiento concurrente de matrículas usando Arc<Mutex<Vec<Matricula>>>.
///
/// Arc (Atomic Reference Counting) permite compartir ownership entre hilos.
/// Mutex garantiza acceso exclusivo al vector de matrículas.
/// El compilador verifica que Matricula implemente Send (puede cruzar hilos).
fn procesar_matriculas_concurrente(
    nuevas_matriculas: Vec<Matricula>,
) -> Vec<Matricula> {

    // Arc<Mutex<>> permite acceso compartido y seguro entre hilos
    let registro = Arc::new(Mutex::new(Vec::<Matricula>::new()));

    let mut handles = Vec::new();

    for matricula in nuevas_matriculas {
        // Clonamos el Arc (no el Mutex): solo se copia el puntero de conteo
        let registro_clone = Arc::clone(&registro);

        let handle = thread::spawn(move || {
            // lock() retorna MutexGuard que se libera automáticamente
            // al salir del bloque (RAII)
            let mut lista = registro_clone.lock().unwrap();

            // Validación simple: evitar matrícula duplicada
            let ya_existe = lista.iter().any(|m| {
                m.estudiante_id == matricula.estudiante_id
                    && m.programa_codigo == matricula.programa_codigo
                    && m.periodo == matricula.periodo
            });

            if !ya_existe {
                println!("[Hilo {:?}] Registrando matrícula: {} en {}",
                         thread::current().id(),
                         matricula.estudiante_id,
                         matricula.programa_codigo);
                lista.push(matricula);
            }
        });
        handles.push(handle);
    }

    // Esperar a que todos los hilos terminen
    for h in handles { h.join().unwrap(); }

    // Extraer el vector del Mutex
    Arc::try_unwrap(registro)
        .expect("Aún existen referencias Arc activas")
        .into_inner()
        .unwrap()
}
\end{lstlisting}

\noindent
La comparación con \java\ resulta instructiva. En \java, el acceso concurrente a la
lista de matrículas requeriría \texttt{synchronized} o \texttt{ReentrantLock}, y un
error de omisión ---olvidar el bloque \texttt{synchronized}--- no es detectado por el
compilador: el error se manifiesta en tiempo de ejecución, de manera no determinística,
bajo carga concurrente elevada. En \rust, si el código intenta acceder al
\texttt{Mutex<>} sin adquirir el \textit{lock}, o intenta compartir un tipo que
no implementa \texttt{Send} entre hilos, el compilador rechaza el código con un
error de tipos. La tabla~\ref{tab:java-rust-concurrencia} sintetiza la diferencia.

\begingroup
\small
\begin{longtable}{@{} L{3.5cm} L{5cm} L{5.5cm} @{}}
	% --- Título y encabezado de la primera página ---
	\caption{Concurrencia en \java\ vs.\ \rust.} \label{tab:java-rust-concurrencia} \\
	\toprule
	\textbf{Aspecto} & \textbf{\java} & \textbf{\rust} \\
	\midrule
	\endfirsthead
	
	% --- Encabezado para las páginas siguientes ---
	\multicolumn{3}{c}{{\bfseries \tablename\ \thetable{} -- Continuación}} \\
	\toprule
	\textbf{Aspecto} & \textbf{\java} & \textbf{\rust} \\
	\midrule
	\endhead
	
	% --- Pie de tabla para páginas intermedias ---
	\midrule
	\multicolumn{3}{r}{{Continúa en la siguiente página...}} \\
	\bottomrule
	\endfoot
	
	% --- Pie de tabla final ---
	\bottomrule
	\endlastfoot
	
	% --- Datos de la tabla ---
	Mecanismo de acceso exclusivo &
	\texttt{synchronized}, \texttt{ReentrantLock} &
	\texttt{Mutex<T>}: el tipo garantiza exclusión por diseño \\
	\midrule
	
	Compartir entre hilos &
	Cualquier objeto es compartible; los errores se detectan en ejecución &
	Solo tipos que implementen \texttt{Send}; verificado por el compilador \\
	\midrule
	
	Detección de \textit{data race} &
	Herramientas externas o revisión manual de código &
	Compilador: imposible compilar código con \textit{data race} (\textit{Safety}) \\
	\midrule
	
	Compartir propiedad &
	Implícito por referencia (punteros gestionados por GC) &
	\texttt{Arc<T>}: conteo de referencias atómico y explícito \\
	\midrule
	
	Costo en ejecución &
	Impacto del GC + overhead de sincronización dinámica &
	Solo el costo mínimo del \texttt{Mutex}, sin pausas de GC \\
\end{longtable}
\endgroup
% ─────────────────────────────────────────────────────────────────────────────
\section{Interfaz gráfica en \rust}
\label{sec:rust-gui}
% ─────────────────────────────────────────────────────────────────────────────

\noindent
El ecosistema de interfaces gráficas en \rust\ es más fragmentado que el de \java.
La biblioteca \texttt{eframe}/\texttt{egui} se destaca por su enfoque de interfaz
inmediata (\textit{immediate mode GUI}), simplicidad de instalación como crate y
capacidad de compilación tanto para escritorio como para WebAssembly. Su modelo de
programación difiere radicalmente de JavaFX: en lugar de definir un árbol de componentes
que persiste entre frames, la interfaz se redibuja completa en cada frame,
y su estado reside en el struct de la aplicación.

\begin{lstlisting}[style=rustStyle, caption={Interfaz gráfica básica con \texttt{eframe}/\texttt{egui}}, label={lst:egui}]
// Cargo.toml (dependencias relevantes):
// [dependencies]
// eframe = "0.27"
// egui   = "0.27"

use eframe::egui;

/// Estado de la aplicación (persiste entre frames).
struct AppEstatutario {
    facultades:      Vec<Facultad>,
    facultad_nueva:  String,
    codigo_nuevo:    String,
    decano_nuevo:    String,
    mensaje_estado:  String,
}

impl AppEstatutario {
    fn new(_cc: &eframe::CreationContext<'_>) -> Self {
        AppEstatutario {
            facultades:     Vec::new(),
            facultad_nueva: String::new(),
            codigo_nuevo:   String::new(),
            decano_nuevo:   String::new(),
            mensaje_estado: String::new(),
        }
    }
}

impl eframe::App for AppEstatutario {

    /// Se invoca en cada frame; redefine la interfaz completa.
    fn update(&mut self, ctx: &egui::Context, _frame: &mut eframe::Frame) {
        egui::TopBottomPanel::top("menu_bar").show(ctx, |ui| {
            ui.heading("Sistema de Información Estatutario — Universidad Distrital");
        });

        egui::SidePanel::left("panel_formulario").show(ctx, |ui| {
            ui.label("Nueva Facultad");
            ui.horizontal(|ui| {
                ui.label("Código:");
                ui.text_edit_singleline(&mut self.codigo_nuevo);
            });
            ui.horizontal(|ui| {
                ui.label("Nombre:");
                ui.text_edit_singleline(&mut self.facultad_nueva);
            });
            ui.horizontal(|ui| {
                ui.label("Decano/a:");
                ui.text_edit_singleline(&mut self.decano_nuevo);
            });
            if ui.button("Agregar").clicked() {
                if !self.codigo_nuevo.is_empty() && !self.facultad_nueva.is_empty() {
                    self.facultades.push(Facultad::new(
                        &self.facultad_nueva,
                        &self.codigo_nuevo,
                        &self.decano_nuevo,
                    ));
                    self.mensaje_estado = format!(
                        "Facultad '{}' agregada.", self.facultad_nueva);
                    self.facultad_nueva.clear();
                    self.codigo_nuevo.clear();
                    self.decano_nuevo.clear();
                } else {
                    self.mensaje_estado =
                        "Complete código y nombre.".to_string();
                }
            }
            ui.label(&self.mensaje_estado);
        });

        egui::CentralPanel::default().show(ctx, |ui| {
            ui.heading("Facultades registradas");
            egui::ScrollArea::vertical().show(ui, |ui| {
                for facultad in &self.facultades {
                    ui.horizontal(|ui| {
                        ui.label(format!("[{}]", facultad.codigo));
                        ui.label(&facultad.nombre);
                        ui.label(format!("— Decano/a: {}", facultad.decano));
                    });
                    ui.separator();
                }
            });
        });
    }
}

fn main() -> Result<(), eframe::Error> {
    let opciones = eframe::NativeOptions::default();
    eframe::run_native(
        "Sistema Estatutario UD",
        opciones,
        Box::new(|cc| Box::new(AppEstatutario::new(cc))),
    )
}
\end{lstlisting}

\noindent
El modelo de interfaz inmediata de \texttt{egui} resulta sorprendente para quienes
provienen de JavaFX o cualquier \textit{framework} basado en árbol de componentes
persistente. No existen \textit{listeners} de eventos ni \textit{callbacks} explícitos:
el método \texttt{update} examina el estado de la interfaz en cada frame y actualiza
el estado de la aplicación en consecuencia. Este modelo simplifica el razonamiento
sobre el estado de la interfaz, toda vez que el estado completo reside en un único
struct que el método \texttt{update} puede leer y modificar sin ambigüedad.

% ─────────────────────────────────────────────────────────────────────────────
\section{Persistencia en \rust}
\label{sec:rust-persistencia}
% ─────────────────────────────────────────────────────────────────────────────

\noindent
La serialización y deserialización en \rust\ se implementa mediante las crates
\texttt{serde} y \texttt{serde\_json}. La integración es especialmente fluida gracias
a la macro \texttt{\#[derive(Serialize, Deserialize)]}, que genera automáticamente
el código de serialización para cualquier struct o enum anotado, sin requerir
escritura manual de conversores ni configuración adicional.

\begin{lstlisting}[style=rustStyle, caption={Persistencia con \texttt{serde} y \texttt{serde\_json}}, label={lst:serde}]
// Cargo.toml:
// [dependencies]
// serde      = { version = "1", features = ["derive"] }
// serde_json = "1"

use serde::{Deserialize, Serialize};
use std::fs;
use std::path::Path;

// Los derives Serialize y Deserialize se heredan en todos los campos
// siempre que estos también los implementen.
#[derive(Debug, Clone, Serialize, Deserialize)]
pub struct Facultad {
    pub nombre:   String,
    pub codigo:   String,
    pub decano:   String,
    pub escuelas: Vec<Escuela>,
}

// ─── Funciones de persistencia ───────────────────────────────────────────────

/// Guarda una lista de Facultades en un archivo JSON.
///
/// # Errors
/// Retorna Err si no se puede escribir el archivo.
pub fn guardar_facultades(
    facultades: &[Facultad],
    ruta: &str,
) -> Result<(), Box<dyn std::error::Error>> {
    let json = serde_json::to_string_pretty(facultades)?;
    // Crear directorios intermedios si no existen
    if let Some(parent) = Path::new(ruta).parent() {
        fs::create_dir_all(parent)?;
    }
    fs::write(ruta, json)?;
    println!("Facultades guardadas en '{}'.", ruta);
    Ok(())
}

/// Carga la lista de Facultades desde un archivo JSON.
///
/// Si el archivo no existe, retorna una lista vacía sin error.
///
/// # Errors
/// Retorna Err si el archivo existe pero no puede ser deserializado.
pub fn cargar_facultades(
    ruta: &str,
) -> Result<Vec<Facultad>, Box<dyn std::error::Error>> {
    if !Path::new(ruta).exists() {
        return Ok(Vec::new());
    }
    let json = fs::read_to_string(ruta)?;
    let facultades: Vec<Facultad> = serde_json::from_str(&json)?;
    Ok(facultades)
}

/// Guarda cualquier tipo serializable en JSON (función genérica).
pub fn guardar_json<T: Serialize>(
    datos: &T,
    ruta: &str,
) -> Result<(), Box<dyn std::error::Error>> {
    let json = serde_json::to_string_pretty(datos)?;
    if let Some(parent) = Path::new(ruta).parent() {
        fs::create_dir_all(parent)?;
    }
    fs::write(ruta, json)?;
    Ok(())
}

/// Carga cualquier tipo deserializable desde JSON (función genérica).
pub fn cargar_json<T: for<'de> Deserialize<'de>>(
    ruta: &str,
) -> Result<Option<T>, Box<dyn std::error::Error>> {
    if !Path::new(ruta).exists() {
        return Ok(None);
    }
    let json = fs::read_to_string(ruta)?;
    let valor: T = serde_json::from_str(&json)?;
    Ok(Some(valor))
}

// ─── Ejemplo de uso ──────────────────────────────────────────────────────────

fn ejemplo_persistencia() -> Result<(), Box<dyn std::error::Error>> {
    let mut fac = Facultad {
        nombre:   "Facultad de Ingeniería".to_string(),
        codigo:   "FI".to_string(),
        decano:   "Dr. Ejemplo".to_string(),
        escuelas: Vec::new(),
    };

    let facultades = vec![fac];
    guardar_facultades(&facultades, "datos/facultades.json")?;

    let cargadas = cargar_facultades("datos/facultades.json")?;
    println!("Facultades cargadas: {}", cargadas.len());
    for f in &cargadas {
        println!("  - {} ({})", f.nombre, f.codigo);
    }
    Ok(())
}
\end{lstlisting}

\noindent
La comparación con la persistencia en \java\ (Gson) revela una diferencia conceptual
significativa. En \java, Gson opera mediante reflexión en tiempo de ejecución: analiza
los campos del objeto en tiempo de ejecución y los serializa dinámicamente. Este enfoque
es flexible pero conlleva un costo: los errores de tipo (por ejemplo, intentar deserializar
un campo \texttt{int} desde un valor JSON de tipo \texttt{string}) se detectan en ejecución,
no en compilación. En \rust, el código de serialización se genera en tiempo de compilación
mediante macros procedurales: los errores de tipo son detectados antes de que el programa
se ejecute, y el rendimiento es superior al no requerir reflexión.

% ─────────────────────────────────────────────────────────────────────────────
\section{\java\ versus \rust: análisis comparativo ampliado}
\label{sec:rust-java-comparacion}
% ─────────────────────────────────────────────────────────────────────────────

\noindent
La implementación del mismo dominio en ambos lenguajes permite una comparación
fundamentada en evidencia concreta, en lugar de los debates abstractos que suelen
dominar las discusiones sobre lenguajes de programación. La tabla~\ref{tab:comparacion-8d}
examina ocho dimensiones relevantes para el sistema estatutario.

\begin{longtable}{@{}L{2.5cm}L{5cm}L{5cm}@{}}
\caption{Comparación de java y rust en ocho dimensiones}
\label{tab:comparacion-8d}\\
\toprule
\textbf{Dimensión} & \textbf{\java} & \textbf{\rust} \\
\midrule
\endfirsthead
\multicolumn{3}{l}{\small\textit{Continuación de la tabla~\ref{tab:comparacion-8d}}}\\
\toprule
\textbf{Dimensión} & \textbf{\java} & \textbf{\rust} \\
\midrule
\endhead
\bottomrule
\endfoot
Herencia &
  Herencia simple de clases (\texttt{extends}); herencia múltiple solo de interfaces.
  Jerarquía \texttt{UnidadAcademica $\leftarrow$ Facultad} &
  No existe herencia de structs. Composición de traits: \texttt{impl EntidadAcademica for Facultad}.
  La composición reemplaza la herencia \\
\addlinespace
Traits vs.\ interfaces &
  Las interfaces definen contratos de comportamiento. Las clases abstractas pueden
  incluir estado. \texttt{default} methods desde Java~8 &
  Los traits definen comportamiento y pueden tener implementaciones por defecto.
  No tienen estado propio. Un struct puede implementar múltiples traits \\
\addlinespace
Null safety &
  \texttt{null} es un valor legal para cualquier referencia de objeto.
  \texttt{NullPointerException} es el error más frecuente en producción.
  Mitigación: \texttt{Optional\textless T\textgreater} (no obligatoria) &
  \texttt{null} no existe. Los valores ausentes se expresan como \texttt{Option\textless T\textgreater}.
  El compilador obliga a tratar \texttt{Some} y \texttt{None}.
  Imposible el \textit{null dereference} \\
\addlinespace
Manejo de errores &
  Excepciones checked y unchecked. La jerarquía \texttt{EstatutoException}
  obliga a declarar y capturar. Riesgo de excepciones no controladas (unchecked) &
  \texttt{Result\textless T, E\textgreater}: el tipo de retorno anuncia explícitamente el posible error.
  El operador \texttt{?} propaga el error. No existen excepciones no controladas \\
\addlinespace
Memoria &
  Recolector de basura (G1GC, ZGC). El programador no gestiona memoria.
  Posibles pausas GC de 1--100 ms en cargas altas &
  Ownership + borrowing. Sin GC. Garantías en tiempo de compilación.
  Rendimiento predecible y sin pausas involuntarias \\
\addlinespace
Concurrencia &
  \texttt{synchronized}, \texttt{Lock}, \texttt{ExecutorService}.
  Las condiciones de carrera se detectan en ejecución (ThreadSanitizer).
  El compilador no garantiza ausencia de \textit{data races} &
  \texttt{Arc\textless Mutex\textless T\textgreater\textgreater}. El sistema de tipos (traits \texttt{Send}/\texttt{Sync})
  garantiza en compilación la ausencia de \textit{data races}.
  Error de concurrencia = error de compilación \\
\addlinespace
Rendimiento &
  JIT: rendimiento excelente tras la fase de calentamiento (warm-up).
  GC introduce latencia variable. Overhead de la JVM en arranque &
  Compilación AOT a código nativo. Rendimiento comparable a C/C++.
  Arranque instantáneo. Sin overhead de VM \\
\addlinespace
Ecosistema &
  Ecosistema maduro: Maven Central, Spring, Hibernate, JUnit.
  Documentación exhaustiva. Comunidad masiva &
  Ecosistema en crecimiento acelerado: crates.io, tokio, serde, egui.
  Curva de aprendizaje pronunciada. Comunidad técnica de alta calidad \\
\end{longtable}

\noindent
Ninguno de los dos lenguajes es superior en términos absolutos. La elección entre
\java\ y \rust\ depende del contexto: un sistema de información universitaria de
gestión académica, con carga moderada y prioridad en la velocidad de desarrollo,
puede estar mejor servido por \java\ y su ecosistema maduro. Un componente de
procesamiento de matrículas en tiempo real, con requisitos estrictos de latencia y
garantías de seguridad de memoria, se beneficiaría de \rust. El presente libro no
propone una respuesta única; propone, en cambio, que el estudiante adquiera
suficiente conocimiento de ambos paradigmas para razonar sobre esa elección con
criterio fundamentado.

% ─────────────────────────────────────────────────────────────────────────────
\section*{Sprints 8 y 9: implementación en \rust}
\label{sec:sprints-8-9}
\addcontentsline{toc}{section}{Sprints 8 y 9: implementación en \rust}
% ─────────────────────────────────────────────────────────────────────────────

\begin{tcolorbox}[sprintBox, title={Sprint 8 — Modelado del dominio en \rust}]
\textbf{Duración:} 2 semanas\\[4pt]
\textbf{Objetivo:} Implementar el modelo del dominio completo en \rust,
con gestión de errores mediante \texttt{Result}/\texttt{Option} y
serialización con \texttt{serde}.

\vspace{6pt}
\begin{tabularx}{\linewidth}{@{}L{4cm}L{4cm}L{4.5cm}@{}}
\toprule
\textbf{Entregable técnico} & \textbf{Característica del programa} & \textbf{Criterio de la rúbrica} \\
\midrule
Módulo \texttt{modelo} con todos los structs y enums del dominio &
  Structs: Facultad, Escuela, Centro, Instituto, Docente, Estudiante, ProgramaAcademico &
  Código compila sin advertencias (\texttt{cargo build}); structs derivados de \texttt{Debug}, \texttt{Clone}, \texttt{Serialize}, \texttt{Deserialize} \\
\addlinespace
Trait \texttt{EntidadAcademica} implementado para Facultad, Escuela, Centro e Instituto &
  Polimorfismo mediante traits; despacho estático y dinámico &
  El trait está correctamente implementado para cada struct; se demuestra uso con \texttt{dyn EntidadAcademica} \\
\addlinespace
Funciones de persistencia \texttt{guardar\_json} y \texttt{cargar\_json} &
  Serialización con \texttt{serde\_json}; archivos en \texttt{datos/} &
  Los datos se persisten correctamente y se recuperan tras reiniciar \\
\addlinespace
Suite de tests (\texttt{\#[cfg(test)]}) con cobertura de casos límite &
  Manejo de errores con \texttt{Result} y \texttt{Option} &
  Todos los tests pasan con \texttt{cargo test}; se prueban casos \texttt{Some}, \texttt{None}, \texttt{Ok}, \texttt{Err} \\
\bottomrule
\end{tabularx}

\vspace{8pt}
\textbf{Definición de terminado:} \texttt{cargo build} y \texttt{cargo test} pasan sin
errores ni advertencias; la serialización de una Facultad con al menos una Escuela y
un Docente produce un archivo JSON válido que puede ser recargado correctamente.
\end{tcolorbox}

\vspace{10pt}

\begin{tcolorbox}[actividadIA, title={Actividad IA — Sprint 8: Ownership y el asistente generativo}]
\textbf{Contexto:} El sistema de ownership de \rust\ es el concepto que más
dificultad genera en estudiantes provenientes de \java\ o Python. Se propone
usar el asistente generativo como interlocutor para aclarar dudas de compilación,
no como generador de soluciones completas.

\vspace{6pt}
\textbf{Tarea:}
\begin{enumerate}
  \item Escribir el siguiente código deliberadamente erróneo y registrar el
  mensaje de error del compilador:
\begin{lstlisting}[style=rustStyle, numbers=none]
let escuela = Escuela::new("EIS", "EIS-001", "Director");
let ref1 = &escuela;
let ref2 = &mut escuela; // ERROR: no se puede tomar
                          // referencia mutable mientras
                          // existe una inmutable
println!("{}", ref1.nombre);
\end{lstlisting}
  \item Pegar el mensaje de error en el asistente generativo y pedirle que lo explique.
  \item Contrastar la explicación del asistente con la sección~\ref{subsec:ownership}.
  \item Identificar si el asistente omitió algún aspecto relevante del modelo de borrowing.
\end{enumerate}

\vspace{6pt}
\textbf{Reflexión:} El compilador de \rust\ genera mensajes de error entre los más
descriptivos y pedagógicos de cualquier lenguaje de programación. Con frecuencia,
el mensaje no solo identifica el error sino que propone la corrección. ¿En qué
medida el asistente generativo añade valor respecto del mensaje del propio compilador?
\end{tcolorbox}

\vspace{10pt}
%\begin{comment}

\begin{tcolorbox}[sprintBox, title={Sprint 9 — Concurrencia, GUI y comparación Java vs.\ Rust}]
	\textbf{Duración:} 2 semanas\\[4pt]
	\textbf{Objetivo:} Implementar procesamiento concurrente de matrículas, una interfaz
	gráfica básica con \texttt{egui} y elaborar el informe comparativo final Java vs.\ Rust.
	
	\vspace{6pt}
	
	\noindent
	\begin{center}
		\begin{tabular}{p{4cm}p{4cm}p{4.5cm}}
			\toprule
			\textbf{Entregable técnico} & \textbf{Característica del programa} & \textbf{Criterio de la rúbrica} \\
			\midrule
			Función \verb|procesar_matriculas_concurrente| con \verb|Arc<Mutex<Matricula>>| &
			Concurrencia segura; evita duplicados en acceso paralelo &
			El código compila sin \verb|unsafe|; se demuestran al menos 4 hilos concurrentes sin \textit{data race} \\
			\addlinespace
			Interfaz gráfica con \texttt{egui}: lista de Facultades y formulario de creación &
			Interfaz inmediata; estado en \texttt{struct} de aplicación &
			La aplicación muestra la lista de facultades persistidas y permite agregar nuevas \\
			\addlinespace
			Informe comparativo Java vs.\ Rust (5 páginas) &
			Análisis fundamentado en la implementación realizada &
			El informe cita evidencia concreta del código propio (no genérico); usa la tabla de 8 dimensiones como estructura \\
			\addlinespace
			Reflexión individual: ownership vs.\ GC (1 página) &
			Perspectiva personal sobre la diferencia de modelo mental &
			La reflexión es específica: hace referencia a al menos un error de ownership que el estudiante cometió y corrigió \\
			\bottomrule
		\end{tabular}
	\end{center}
	
	\vspace{8pt}
	\textbf{Definición de terminado:} \verb|cargo run| abre la ventana de \texttt{egui};
	el test de concurrencia procesa 20 matrículas en paralelo y el resultado no contiene
	duplicados; el informe comparativo supera la revisión de pares del equipo.
\end{tcolorbox}

%\end{comment}

\vspace{10pt}

\begin{tcolorbox}[actividadIA, title={Actividad IA — Sprint 9: Comparación asistida \java\ vs.\ \rust}]
\textbf{Tarea:}
\begin{enumerate}
  \item Solicitar al asistente generativo: \textit{``Compara el manejo de
  concurrencia en Java y Rust para el caso específico de procesamiento paralelo
  de matrículas universitarias. Enfócate en las garantías de seguridad que
  ofrece cada lenguaje y en los mecanismos de sincronización disponibles.''}
  \item Contrastar la respuesta del asistente con la tabla~\ref{tab:java-rust-concurrencia}
  y el código de los listados~\ref{lst:concurrencia}.
  \item Identificar afirmaciones del asistente que no se verifican en la
  implementación concreta realizada en los sprints 6--9.
  \item Calcular el porcentaje de afirmaciones del asistente que son
  correctas, parcialmente correctas y erróneas, documentando la metodología.
\end{enumerate}

\vspace{6pt}
\textbf{Punto de reflexión metodológica:} La evaluación crítica de la calidad
de una respuesta generada por IA requiere conocimiento previo del dominio.
Un estudiante que no ha implementado concurrencia en ambos lenguajes no puede
evaluar si la respuesta del asistente es correcta. Esto no invalida el uso
de asistentes generativos; redefine su lugar en el proceso de aprendizaje:
son más útiles para el experto que para el principiante, aun cuando el
principiante sea quien más los use.
\end{tcolorbox}
